\section{Tres niveles de acceso: API, Open Weights, Open Source}

\subsection{El modelo de tres niveles de acceso}
El expositor divide el acceso en tres niveles, cada uno correspondiente a una posición de investigación distinta:
\begin{itemize}
    \item API access: como un «científico conductual» que observa las entradas y salidas de una caja negra;
    \item Open weights: como un «neurocientífico» que puede examinar las representaciones internas;
    \item Open source: como un «ingeniero de sistemas» que puede reescribir todo el pipeline.
\end{itemize}

\begin{table}[H]
\centering
\begin{tabular}{p{2.4cm}p{4.2cm}p{4.8cm}}
\toprule
\textbf{Nivel de acceso} & \textbf{Capacidades típicas} & \textbf{Limitaciones principales} \\
\midrule
API access & Diseño de prompts, orquestación de llamadas a herramientas, evaluación de caja negra & Difícil verificar las causas mecánicas, difícil hacer modificaciones a nivel estructural \\
Open weights & Análisis de representaciones, destilación/fine-tuning, experimentos de intervención & Aun así, sigue restringido por la receta de entrenamiento y la arquitectura existentes \\
Open source & Reconstrucción de toda la cadena de datos, modelo, entrenamiento y evaluación & Exige mucho de la ingeniería y del cómputo, el costo de reproducción es alto \\
\bottomrule
\end{tabular}
\caption{El marco de «niveles de acceso» propuesto por el expositor y sus fronteras de investigación}
\end{table}

\begin{importantbox}{Un mismo modelo, distintos niveles de acceso llevan a distintos tipos de artículo}
El recordatorio más valioso del curso es este: la pregunta de investigación no está determinada únicamente por el interés, sino también por «qué experimento es posible hacer».
En la era de la API son más frecuentes los artículos de estrategia o de nivel de prompt; en la era de open source son más probables los artículos de arquitectura, objetivo de entrenamiento, construcción de datos y codiseño de sistemas.
\end{importantbox}

\subsection{Por qué hay que distinguir Open Weights de Open Source}
\begin{knowledgebox}{Open Weights no equivale a Open Science}
Tener los pesos permite hacer muchas cosas, pero aun así no se sabe:
\begin{itemize}
    \item cómo influye la proporción de mezcla de datos en las fronteras de las capacidades;
    \item cómo cambia el comportamiento el currículo de entrenamiento (curriculum) y las estrategias de alineación;
    \item si la mejora de desempeño proviene de la arquitectura o de la ingeniería de datos.
\end{itemize}
Por eso, Open Weights es un estado intermedio importante, pero todavía no llega al estado científico de «poder verificarse por completo».
\end{knowledgebox}

\begin{warningbox}{Riesgo de confusión conceptual en las discusiones de política}
En las discusiones de la industria suele mezclarse «poder invocar la API», «poder descargar los pesos» y «poder reproducir el entrenamiento» en una sola palabra: «abierto». Esto hace que el diseño de la regulación y de los mecanismos de colaboración se aleje de las restricciones técnicas reales y genere incentivos equivocados.
\end{warningbox}

\subsection{Resumen de la sección}
El marco de tres niveles que propone el expositor permite convertir el «debate sobre la apertura» de un eslogan en un problema de ingeniería: distintos niveles corresponden a distintas capacidades experimentales y distintas exposiciones al riesgo, y no deben confundirse entre sí.

